Soit $n \geqslant 1$ un entier. Un \emph{triangle nippon} consiste en $1 + 2 + \cdots + n$ cercles disposés
de manière à constituer un triangle formé de $n$ lignes, dont la $\ell$\textsuperscript{e} ligne contient
exactement $\ell$ cercles et a été coloriée de sorte qu'un de ces $\ell$ cercles soit rouge et que les autres
cercles de la ligne soient incolores. On appelle \emph{chemin ninja} toute suite de $n$ cercles qui débute
sur le cercle de la première ligne, dont chaque cercle est situé juste en-dessous du cercle qui le précède,
et qui se termine sur un cercle de la dernière ligne.

(L'énoncé officiel donne ici, en figure, un exemple de triangle nippon pour $n = 6$, avec un chemin ninja
contenant deux cercles rouges.)

Trouver, en fonction de $n$, le plus grand entier $k$ pour lequel, dans chaque triangle nippon, il existe
nécessairement un chemin ninja contenant au moins $k$ cercles rouges.
